% Shared preamble.  Kept separate so the document body stays readable.
\usepackage[UTF8]{ctex}
\setCJKmainfont{Noto Serif CJK SC}
\setCJKsansfont{Noto Sans CJK SC}
\setCJKmonofont{Noto Sans Mono CJK SC}
\usepackage{amsmath,amssymb,amsthm}
\usepackage{booktabs,multirow,array}
\usepackage{algorithm}
\usepackage[noend]{algpseudocode}
\usepackage{graphicx}
\graphicspath{{../../results/figures_r14/}}
\usepackage[table]{xcolor}
\usepackage{geometry}
\geometry{a4paper,margin=2.4cm}
\usepackage[colorlinks=true,linkcolor=blue!55!black,citecolor=blue!55!black,
            urlcolor=blue!55!black]{hyperref}
\usepackage{caption}
\captionsetup{font=small,labelfont=bf}
\usepackage{enumitem}
\setlist{itemsep=2pt,parsep=0pt,topsep=3pt}
\usepackage{fancyhdr}
\setlength{\headheight}{14pt}
\addtolength{\topmargin}{-2pt}
\pagestyle{fancy}\fancyhf{}
\fancyhead[L]{\small 几何感知张量分解}
\fancyhead[R]{\small\thepage}
\renewcommand{\headrulewidth}{0.4pt}

% Algorithm environment in Chinese.
\floatname{algorithm}{算法}
\renewcommand{\algorithmicrequire}{\textbf{输入：}}
\renewcommand{\algorithmicensure}{\textbf{输出：}}
\algrenewcommand\algorithmicfor{\textbf{对}}
\algrenewcommand\algorithmicdo{\textbf{执行}}
\algrenewcommand\algorithmicwhile{\textbf{当}}
\algrenewcommand\algorithmicif{\textbf{若}}
\algrenewcommand\algorithmicthen{\textbf{则}}
\algrenewcommand\algorithmicelse{\textbf{否则}}
\algrenewcommand\algorithmicend{\textbf{结束}}
\algrenewcommand\algorithmicreturn{\textbf{返回}}

% Notation shorthands used throughout.
\newcommand{\bY}{\mathcal{Y}}
\newcommand{\bG}{\mathcal{G}}
\newcommand{\bA}{\mathcal{A}}
\newcommand{\Rb}{\mathbb{R}}
\newcommand{\Om}{\Omega}
\newcommand{\Sen}{\mathcal{S}}
\newcommand{\norm}[1]{\lVert #1 \rVert}
\newcommand{\code}[1]{\texttt{\small #1}}

\theoremstyle{definition}
\newtheorem{principle}{准则}
